\usepackage[paperwidth=170mm,paperheight=240mm,inner=22mm,outer=19mm,top=23mm,bottom=24mm,headheight=15pt]{geometry}
\usepackage{amsmath,amssymb,amsthm,mathtools}
\usepackage{booktabs,longtable,array,tabularx}
\usepackage{graphicx,xcolor,tikz}
\usetikzlibrary{arrows.meta,positioning,calc,fit,backgrounds}
\usepackage{enumitem,fancyhdr,microtype}
\usepackage[numbers,sort&compress]{natbib}
\usepackage{hyperref}
\definecolor{ink}{HTML}{17334B}
\definecolor{accent}{HTML}{A45335}
\definecolor{pale}{HTML}{E9EFF3}
\hypersetup{unicode=true,colorlinks=true,linkcolor=ink,citecolor=accent,urlcolor=ink,pdftitle={思维论},pdfsubject={诊断组织的有限数学：共同错误、可辨别性与通信},pdfauthor={AI员工项目研究书稿}}
\urlstyle{same}
\setlength{\emergencystretch}{2em}
\setlength{\parskip}{0.25em}
\setlength{\tabcolsep}{5pt}
\linespread{1.12}
\raggedbottom
\setcounter{tocdepth}{1}
\setcounter{secnumdepth}{2}
\setlist{nosep,topsep=4pt,leftmargin=2em}
\pagestyle{fancy}
\fancyhf{}
\fancyhead[LE,RO]{\small\thepage}
\fancyhead[RE]{\small 思维论}
\fancyhead[LO]{\small\nouppercase{\leftmark}}
\renewcommand{\headrulewidth}{0.3pt}
\ctexset{chapter={format=\Large\bfseries\color{ink},beforeskip=12pt,afterskip=22pt},section={format=\large\bfseries\color{ink}}}
\theoremstyle{definition}
\newtheorem{definition}{定义}[chapter]
\newtheorem{example}[definition]{例}
\newtheorem{counterexample}[definition]{反例}
\newtheorem{exercise}{练习}[chapter]
\theoremstyle{plain}
\newtheorem{theorem}[definition]{定理}
\newtheorem{proposition}[definition]{命题}
\newtheorem{lemma}[definition]{引理}
\newtheorem{corollary}[definition]{推论}
\theoremstyle{remark}
\newtheorem{remark}[definition]{说明}
\renewcommand{\proofname}{证明}
\newcommand{\VV}{\mathcal V}
\newcommand{\HH}{\mathcal H}
\newcommand{\LL}{\mathcal L}
\newcommand{\XX}{\mathcal X}
\newcommand{\II}{\mathbf I}
\newcommand{\PP}{\mathbb P}
\newcommand{\EE}{\mathbb E}
\newcommand{\NN}{\mathbb N}
\newcommand{\RR}{\mathbb R}
\newcommand{\concat}{\mathbin{\Vert}}
\newcommand{\wait}{\mathsf{wait}}
\newcommand{\StopMode}{\mathsf{stop}}
\newcommand{\TV}{\operatorname{TV}}
\newcommand{\Var}{\operatorname{Var}}
\newcommand{\Cov}{\operatorname{Cov}}
\newcommand{\KL}{D_{\mathrm{KL}}}
\newcommand{\bits}{\operatorname{bits}}
\newcommand{\cost}{\operatorname{cost}}
\newcommand{\obs}{\operatorname{obs}}
\newcommand{\sem}[1]{[\![#1]\!]}
\newcommand{\chapterguide}[1]{\begin{quote}\small\color{ink}#1\end{quote}}
\DeclareMathOperator{\Risk}{Risk}
\DeclareMathOperator{\Reg}{Reg}
\DeclareMathOperator{\Bin}{Bin}
\DeclareMathOperator{\Bern}{Bern}
\newcommand{\calP}{\mathfrak P}
\newcommand{\calD}{\mathcal D}
\newcommand{\calA}{\mathcal A}
\newcommand{\calQ}{\mathcal Q}
\newcommand{\kap}{\kappa_{\mathrm H}}
\newcommand{\kapK}{\kappa_{\mathrm K}}
\newcommand{\Aff}{\operatorname{Aff}}
\newcommand{\Alt}{\operatorname{Alt}}
\newcommand{\kl}{\operatorname{kl}}
\newcolumntype{Y}{>{\raggedright\arraybackslash}X}
\tikzset{fragment/.style={circle,draw=ink,fill=pale,minimum size=7mm,inner sep=1pt},operation/.style={rectangle,draw=accent,fill=white,minimum width=9mm,minimum height=6mm},unit/.style={rounded corners,draw=ink,fill=pale,minimum width=22mm,minimum height=11mm,align=center},edge/.style={-{Stealth[length=2mm]},semithick,draw=ink}}
